% !TEX TS-program = pLaTeX (ptex2pdf)
% Public example: fictitious names, dates and measurements.
\RequirePackage{plautopatch}
\documentclass[platex,dvipdfmx,titlepage,a4j]{jsarticle}

\usepackage[T1]{fontenc}
\usepackage{lmodern}
\usepackage{amsmath,amssymb}
\usepackage{graphicx}
\usepackage[yen]{okuverb}
\usepackage{school-report-example}

\title{抵抗の電圧と電流の関係}
\grade{3年00番}
\author{サンプル 太郎}
\team{第1班}
\date{2026年10月18日}
\expdate{2026年10月15日}
\coauthor{
  01番 & サンプル 花子 \\
  02番 & サンプル 次郎
}

\begin{document}
\maketitle

\section{目的}
抵抗に加える電圧と流れる電流の関係を調べ，
オームの法則と比較する。
本書の氏名・日付・数値は使用例であり，実測値ではない。

\section{理論}
抵抗値を \(R\)，電圧を \(V\)，電流を \(I\) とすると，
オームの法則は式\eqref{eq:ohm}で表される。
\begin{equation}
  I = \frac{V}{R}
  \label{eq:ohm}
\end{equation}

\section{実験方法}
抵抗値を \(1.0\,\mathrm{k\Omega}\) とした回路を考え，
印加電圧を変えたときの電流を比較する。
実際のレポートでは，使用器具，回路図，操作手順を記載する。

\section{実験結果}
説明用の計算値を表\ref{tab:values}に示す。
表番号は \verb|\label| と \verb|\ref| を使って参照する。
\begin{table}[htbp]
  \centering
  \caption{電圧と電流の計算値（仮想データ）}
  \label{tab:values}
  \begin{tabular}{rr}
    \hline
    電圧 / \(\mathrm{V}\) & 電流 / \(\mathrm{mA}\) \\
    \hline
    1.0 & 1.0 \\
    2.0 & 2.0 \\
    3.0 & 3.0 \\
    \hline
  \end{tabular}
\end{table}

\section{考察}
表\ref{tab:values}では，電圧が2倍になると電流も2倍になり，
式\eqref{eq:ohm}の比例関係を確認できる。
実測値を扱う場合は，誤差の大きさとその原因も検討する。

\begin{thebibliography}{9}
  \bibitem{guide}
    サンプル著者，「実験レポートの記述例」，2026（架空の資料）．
\end{thebibliography}
参考文献欄は記入例である。
提出時には，実際に参照した資料を記載する。

\section*{感想}
数式や表を相互参照で記述すると，
内容を変更した際も番号を自動的に更新できることを確認した。

\end{document}
